%
% Main document
% =======================================================================
% This is part of the document "Project documentation template".
% Authors: lk
%----------------------------------------------------------------------
\documentclass[
	a4paper,				% paper format
	10pt,					% fontsize
	%twoside,				% double-sided
	openright,				% begin new chapter on right side
	notitlepage,			% use no standard title page
	parskip=half,			% set paragraph skip to half of a line
]{scrreprt}					% KOMA-script report
%----------------------------------------------------------------------

\raggedbottom
\KOMAoptions{cleardoublepage=plain}		% Add header and footer on blank pages


% Load Standard Packages:
%----------------------------------------------------------------------
\usepackage[standard-baselineskips]{cmbright}

\usepackage[ngerman]{babel}				% german hyphenation
%\usepackage[latin1]{inputenc} 			% Unix/Linux - load extended character set (ISO 8859-1)
\usepackage[ansinew]{inputenc}  		% Windows - load extended character set (ISO 8859-1)
\usepackage[T1]{fontenc}				% hyphenation of words with �,� and �
\usepackage{textcomp}					% additional symbols
\usepackage{ae}							% better resolution of Type1-Fonts 
\usepackage{fancyhdr}					% simple manipulation of header and footer 
\usepackage{etoolbox}					% color manipulation of header and footer
\usepackage{graphicx}                  	% integration of images
\usepackage{float}						% floating objects
\usepackage{caption}					% for captions of figures and tables
\usepackage{booktabs}					% package for nicer tables
\usepackage{tocvsec2}					% provides means of controlling the sectional numbering
\usepackage{lastpage}  					% to include max pagenumber
%----------------------------------------------------------------------

% Load Math Packages
%----------------------------------------------------------------------
\usepackage{amsmath}					% various features to facilitate writing math formulas
\usepackage{amsthm}						% enhanced version of latex's newtheorem
\usepackage{amsfonts}					% set of miscellaneous TeX fonts that augment the standard CM
\usepackage{amssymb}					% mathematical special characters
\usepackage{exscale}					% mathematical size corresponds to textsize
\usepackage{siunitx}					% si units
%----------------------------------------------------------------------

% Package to facilitate placement of boxes at absolute positions
%----------------------------------------------------------------------
\usepackage[absolute]{textpos}
\setlength{\TPHorizModule}{1mm}
\setlength{\TPVertModule}{1mm}
%----------------------------------------------------------------------					
			
% Definition of Colors
%----------------------------------------------------------------------
\RequirePackage{color}							% Color (not xcolor!)
\definecolor{linkblue}{rgb}{0,0,0.8}			% Standard
\definecolor{fhnwyellow}{rgb}{0.253,0.229,0.77}	% FHNW Yellow
\definecolor{linkcolor}{rgb}{0,0,0.8}			% Blue for the web- and cd-version!
\definecolor{linkcolor}{rgb}{0,0,0}				% Black for the print-version!
%----------------------------------------------------------------------

% Hyperref Package (Create links in a pdf)
%----------------------------------------------------------------------
\usepackage[
	pdftex,ngerman,bookmarks,plainpages=false,pdfpagelabels,
	backref = {false},							% No index backreference
	colorlinks = {true},						% Color links in a PDF
	hypertexnames = {true},						% no failures "same page(i)"
	bookmarksopen = {true},						% opens the bar on the left side
	bookmarksopenlevel = {0},					% depth of opened bookmarks
	pdftitle = {Template FHNW},					% PDF-property
	pdfauthor = {lk},							% PDF-property
	pdfsubject = {LaTeX Template},				% PDF-property
	linkcolor = {linkcolor},					% Color of Links
	citecolor = {linkcolor},					% Color of Cite-Links
	urlcolor = {linkcolor},						% Color of URLs
]{hyperref}
%----------------------------------------------------------------------

% Set up page dimension
%----------------------------------------------------------------------
\usepackage{geometry}
\geometry{
	a4paper,
	left=28mm,
	right=15mm,
	top=30mm,
	headheight=20mm,
	headsep=8mm,
	textheight=242mm,
	footskip=15mm
}
%----------------------------------------------------------------------

% Makeindex Package
%----------------------------------------------------------------------
\usepackage{makeidx}							% To produce index
\makeindex										% Index-Initialisation
%----------------------------------------------------------------------

% Glossary Package
%----------------------------------------------------------------------
% the glossaries package uses makeindex if you use TeXnicCenter do the following steps:
% - Goto "Ausgabeprofile definieren" (ctrl + F7)
% - Select the profile "LaTeX => PDF"
% - Add in register "Nachbearbeitung" a new "Postprozessoren" point named Glossar
% - Select makeindex.exe in the field "Anwendung" ( ..\MiKTeX x.x\miktex\bin\makeindex.exe )
% - Add this [ -s "%tm.ist" -t "%tm.glg" -o "%tm.gls" "%tm.glo" ] in the field "Argumente"
%
% for futher informations go to http://ewus.de/tipp-1029.html
%----------------------------------------------------------------------
%\usepackage[nonumberlist]{glossaries}
%\makeglossaries
%
\newglossaryentry{BibTeX}{name={BibTeX},description={Programm zur Erstellung von Literaturangaben und -verzeichnissen in \TeX- oder \LaTeX-Dokumenten}}
\newglossaryentry{StwVrz}{name={Stichwortverzeichnis},description={Verzeichnis mit Stichworten aus dem Text}}



%----------------------------------------------------------------------

% Intro:
%----------------------------------------------------------------------
\begin{document}							% Start Document
\settocdepth{section}						% Set depth of toc
\pagenumbering{roman}		

% Report definitions
%----------------------------------------------------------------------
\renewcommand{\tablename}{Tableau}			% Tag for Figures
\renewcommand{\figurename}{Fig.}			% Tag for Figures
\providecommand{\titel}{Titel des Berichts}			%  Hier den Titel des Berichts/Thesis eingeben
\providecommand{\untertitel}{Untertitel}			%  Untertitel des Berichts
\providecommand{\berichtklasse}{Weitere Titel}		%  z.B. Laborbericht
\providecommand{\studiengang}{Bauingenieurwesen}	%  Studiengang
\providecommand{\autoren}{Vorname, Name}			%  Autoren
\providecommand{\betreuer}{Prof. Dr.~X}				%  betreuender Dozent
\providecommand{\versionnumber}{1.0}				%  Hier die aktuelle Versionsnummer eingeben
\providecommand{\versiondate}{\today}				%  Standarddatum: Tag der letzten Kompilation\textsl				% Titel der Arbeit aus Datei titel.tex lesen


% Set up header and footer
%----------------------------------------------------------------------
\makeatletter
\patchcmd{\@fancyhead}{\rlap}{\color{black}\rlap}{}{}		% new color of header
\patchcmd{\@fancyfoot}{\rlap}{\color{black}\rlap}{}{}		% new color of footer
\makeatother

\fancyhf{}													% clean all fields
\fancypagestyle{plain}{										% new definition of plain style	
	\fancyfoot[R]{\footnotesize \thepage} 					% footer right part --> page number
	\fancyfoot[L]{\footnotesize \titel, Version \versionnumber, \versiondate}											% footer even page left part 
	\lhead{\includegraphics[scale=0.8]{bilder/fhnw_logo.png}}																% set logo on each page
}
\fancyfoot[R]{\normalsize \thepage}
\renewcommand{\chaptermark}[1]{\markboth{\thechapter.  #1}{}}
\renewcommand{\headrulewidth}{2pt}				% no header stripline
\renewcommand{\footrulewidth}{2pt} 				% no bottom stripline

\pagestyle{plain}
%----------------------------------------------------------------------


% Title Page and Abstract
%----------------------------------------------------------------------
%\include{vorspann/titelseite_ohne_bild}		% activate for Titelseite ohne Bild
%
% Project documentation template
% =======================================================================
% This is part of the document "Project documentation template".
% Authors: lk
%

\begin{titlepage}

% FHNW-Logo absolute placed at (28,12) on A4 and picture
%----------------------------------------------------------------------
\setlength{\unitlength}{1mm}
\begin{textblock}{20}[0,0](28,12)
	\includegraphics[scale=1.0]{bilder/fhnw_logo.png}
\end{textblock}

\begin{textblock}{154}[0,0](28,50)
	\includegraphics[scale=1.0]{bilder/platzhalter.jpg}			% Titelbild definieren
\end{textblock}

% Institution / Titel / Untertitel / Autoren / Experten:
%----------------------------------------------------------------------
\begin{flushleft}

\vspace*{115mm}

\fontsize{26pt}{28pt}\selectfont 
\titel \\								% Titel aus der Datei vorspann/titel.tex lesen
\vspace{2mm}

%\fontsize{16pt}{20pt}\selectfont\vspace{0.3em}
%\untertitel	\\						% Untertitel eingeben
%\vspace{5mm}

%\fontsize{10pt}{12pt}\selectfont
%\textbf{\berichtklasse} \\				% eingeben
%\vspace{3mm}

% Abstract (eingeben):
%----------------------------------------------------------------------
\begin{textblock}{150}(28,190)
\fontsize{10pt}{12pt}\selectfont		%[Kurztext (Abstract) einf�gen, falls gew�nscht] \\ 
\end{textblock}

\begin{textblock}{150}(28,225)
\fontsize{10pt}{17pt}\selectfont
\begin{tabbing}
xxxxxxxxxxxxxx\=\kill
Studiengang: \> \studiengang \\			% Namen eingeben
Autoren: \> \autoren \\					% Namen eingeben
Betreuer: \> \betreuer \\				% Namen eingeben
%Auftraggeber: \> [X AG] \\				% Namen eingeben
%Experten: \> [Prof. Dr.~X]	\\			% Namen eingeben
Datum: \> \versiondate \\				% aus Datei vorspann/version.tex lesen
\end{tabbing}

\end{textblock}
\end{flushleft}

\begin{textblock}{150}(28,280)
\noindent 
\color{black}\fontsize{9pt}{10pt}\selectfont
Fachhochschule Nordwestschweit | Hochschule f�r Architektur, Bau und Geomatik
\color{black}\selectfont
\end{textblock}

\end{titlepage}

% =======================================================================
%			% activate for Titelseite mit Bild
\cleardoubleemptypage
\setcounter{page}{1}
\cleardoublepage
\phantomsection 
%\addcontentsline{toc}{chapter}{Management Summary}
%\include{vorspann/management_summary}
%\cleardoubleemptypage
%----------------------------------------------------------------------

% Table of contents
%----------------------------------------------------------------------
\tableofcontents
\cleardoublepage
\clearpage
\label{EndContents}
%----------------------------------------------------------------------

% Main part:
%----------------------------------------------------------------------
\pagenumbering{arabic}

\chapter{Kapitel}
\label{chap:kapitel}

Text		% Hier alle Kapitel einf�gen
	
%----------------------------------------------------------------------

% Selbst�ndigkeitserkl�rung
%----------------------------------------------------------------------
%\cleardoublepage
%\phantomsection 
%\addcontentsline{toc}{chapter}{Selbst�ndigkeitserkl�rung}
%\chapter*{Selbst�ndigkeitserkl�rung}
\label{chap:selbstaendigkeitserklaerung}

\vspace*{10mm} 

Ich/wir best�tige/n, dass ich/wir die vorliegende Arbeit selbstst�ndig und ohne Benutzung anderer als der im Literaturverzeichnis angegebenen Quellen und Hilfsmittel angefertigt habe/n. S�mtliche Textstellen, die nicht von mir/uns stammen, sind als Zitate gekennzeichnet und mit dem genauen Hinweis auf ihre Herkunft versehen. 

\vspace{15mm}

\begin{tabbing}
xxxxxxxxxxxxxxxxxxxxxxxxx\=xxxxxxxxxxxxxxxxxxxxxxxxxxxxxx\=xxxxxxxxxxxxxxxxxxxxxxxxxxxxxx\kill
Ort, Datum: \> Basel, \versiondate \\
\\ 
Namen Vornamen: \> Vorname, Name \> Vorname, Name  \\
\\
\\
\\ 
Unterschriften: \> ......................................\> ...................................... \\
\end{tabbing}

%----------------------------------------------------------------------

% Glossary
%----------------------------------------------------------------------
%\cleardoublepage
%\phantomsection 
%\addcontentsline{toc}{chapter}{Glossar}
%\renewcommand{\glossaryname}{Glossar}
%\printglossary
%----------------------------------------------------------------------

% Bibliography
%----------------------------------------------------------------------
%\cleardoublepage
%\phantomsection 
%\addcontentsline{toc}{chapter}{Literaturverzeichnis}
%\bibliographystyle{IEEEtranS}
%\bibliography{datenbanken/bibliography}{}
%----------------------------------------------------------------------

% Listings
%----------------------------------------------------------------------
%\cleardoublepage
\phantomsection 
\addcontentsline{toc}{chapter}{Abbildungsverzeichnis}
\listoffigures
%\cleardoublepage
\phantomsection 
\addcontentsline{toc}{chapter}{Tabellenverzeichnis}
\listoftables
%----------------------------------------------------------------------

% Index
%----------------------------------------------------------------------
%\cleardoublepage
%\phantomsection 
%\addcontentsline{toc}{chapter}{Stichwortverzeichnis}
%\renewcommand{\indexname}{Stichwortverzeichnis}
%\printindex
%----------------------------------------------------------------------

% Attachment:
%----------------------------------------------------------------------
%\appendix
%\settocdepth{section}
%\chapter{Beliebiger Anhang}
\label{chap:bel_anhang}

Text

%----------------------------------------------------------------------

%----------------------------------------------------------------------
\end{document}

